\subsection{Formalization}
\label{sec:formalization}

This section presents formal mathematical definitions for the three proposed inconsistency measures. The formalization builds upon the planning foundations established in Section~\ref{sec:foundations_planning} and provides precise algorithmic descriptions for computing each measure.

\subsubsection{Preliminary Definitions}

Before defining the measures themselves, we establish several foundational concepts required for their formalization.

\begin{defn}[Reachable States]
    Let $\langle P, O, I, G \rangle$ be a planning problem. The set of \emph{reachable states} $S_{\text{reach}}$ is defined as the smallest set satisfying:
    \begin{enumerate}
        \item $I \in S_{\text{reach}}$
        \item If $s \in S_{\text{reach}}$ and $o \in O$ with $\text{precond}(o) \subseteq s$, then $\text{RESULT}(s,o) \in S_{\text{reach}}$
    \end{enumerate}
    where $\text{RESULT}(s,o) = (s \setminus \text{delete}(o)) \cup \text{add}(o)$.
\end{defn}

The set $S_{\text{reach}}$ can be computed via forward reachability analysis: starting from $I$, iteratively apply all applicable operators until no new states are generated (a fixpoint is reached). Since $P$ is finite and each state is a subset of $P$, this process terminates in at most $2^{|P|}$ iterations, though practical implementations typically converge much faster.

\begin{defn}[Forward Reachability from State]
    Let $\langle P, O, I, G \rangle$ be a planning problem. For any state $s \subseteq P$,
    the set of \emph{forward-reachable states from $s$}, denoted $\text{Reach}(s)$, is
    the smallest set satisfying:
    \begin{enumerate}
        \item $s \in \text{Reach}(s)$
        \item If $s' \in \text{Reach}(s)$ and $o \in O$ with $\text{precond}(o) \subseteq s'$,
              then $\text{RESULT}(s',o) \in \text{Reach}(s)$
    \end{enumerate}
\end{defn}

Note that $S_{\text{reach}} = \text{Reach}(I)$, and for any $s \in S_{\text{reach}}$,
we have $\text{Reach}(s) \subseteq S_{\text{reach}}$.

\subsubsection{Proposal 1: Unreachability Profile (Formal)}

The Unreachability Profile quantifies Category~2 inconsistency by analyzing which propositions fail to appear in any reachable state. This profile consists of two complementary measures.

\begin{defn}[Unreachable Goal Propositions Count]
    Let $\langle P, O, I, G \rangle$ be a planning problem with reachable states $S_{\text{reach}}$. The measure $I^{\text{scope}}_{\text{UR}}$ counting goal propositions that are unreachable is defined as:
    \begin{equation*}
        I^{\text{scope}}_{\text{UR}}(\langle P, O, I, G \rangle) = |\{g \in G \mid \nexists s \in S_{\text{reach}} : g \in s\}|
    \end{equation*}
\end{defn}

This measure indicates the \emph{scope} of unreachability problems at the goal level: how many goal propositions cannot be achieved. The value ranges from $0$ (all goals reachable) to $|G|$ (no goals reachable).

\begin{defn}[Total Unreachable Propositions Count]
    Let $\langle P, O, I, G \rangle$ be a planning problem with reachable states $S_{\text{reach}}$. The measure $I^{\text{struct}}_{\text{UR}}$ counting all unreachable propositions is defined as:
    \begin{equation*}
        I^{\text{struct}}_{\text{UR}}(\langle P, O, I, G \rangle) = |\{p \in P \mid \nexists s \in S_{\text{reach}} : p \in s\}|
    \end{equation*}
\end{defn}

This measure reveals the \emph{structure} of dependency problems by counting all unreachable propositions throughout the proposition space. The value ranges from $0$ to $|P|$. The difference $I^{\text{struct}}_{\text{UR}} - I^{\text{scope}}_{\text{UR}}$ indicates the depth of cascading dependencies: propositions that are unreachable due to blocked preconditions. Together, $I^{\text{scope}}_{\text{UR}}$ and $I^{\text{struct}}_{\text{UR}}$ form the Unreachability Profile.

\paragraph{Example: Locked Door (Scenario~1).} For the Locked Door scenario defined in Section~\ref{sec:test_scenarios}, forward reachability analysis yields $S_{\text{reach}} = \{\{\}\}$ since no operator is applicable from the initial state. The Unreachability Profile yields:
\begin{align*}
    I^{\text{scope}}_{\text{UR}}  & = |\{\textit{in\_room}\}| = 1                                             \\
    I^{\text{struct}}_{\text{UR}} & = |\{\textit{has\_key}, \textit{door\_unlocked}, \textit{in\_room}\}| = 3
\end{align*}
The difference of 2 indicates two intermediate propositions in the dependency chain that are unreachable due to the missing \textit{has\_key}.

\paragraph{Example: Bank Vault (Scenario~2).} For the Bank Vault scenario, reachability analysis yields $S_{\text{reach}} = \{\{\}\}$ since none of the root preconditions (\textit{key1}, \textit{key2}, \textit{admin\_approval}) are initially true or produced by any operator. The Unreachability Profile yields:
\begin{align*}
    I^{\text{scope}}_{\text{UR}}  & = |\{\textit{in\_vault}\}| = 1                                              \\
    I^{\text{struct}}_{\text{UR}} & = |\{\textit{key1}, \textit{key2}, \textit{admin\_approval},                \\
                                  & \quad \textit{door1\_open}, \textit{door2\_open}, \textit{access\_granted}, \\
                                  & \quad \textit{in\_vault}\}| = 7
\end{align*}

This demonstrates how $I^{\text{struct}}_{\text{UR}}$ distinguishes shallow dependency problems (difference of 2 in Locked Door) from deep cascading failures (difference of 6 in Bank Vault), revealing three independent missing precondition chains converging on a single goal.

\subsubsection{Proposal 2: Goal Mutex Profile (Formal)}

The Goal Mutex Profile quantifies operator-induced mutual exclusions by analyzing which goal pairs can each be achieved but never co-occur in any reachable state. This profile consists of two complementary measures.

\begin{defn}[Mutex Goal Pair]
    \label{defn:mutex}
    Let $\langle P, O, I, G \rangle$ be a planning problem with reachable states $S_{\text{reach}}$. Goal propositions $g_1, g_2 \in G$ with $g_1 \neq g_2$ are \emph{mutex} if:
    \begin{enumerate}
        \item $\exists s_1 \in S_{\text{reach}} : g_1 \in s_1$ (goal $g_1$ is achievable)
        \item $\exists s_2 \in S_{\text{reach}} : g_2 \in s_2$ (goal $g_2$ is achievable)
        \item $\nexists s \in S_{\text{reach}} : \{g_1, g_2\} \subseteq s$ (goals never co-occur)
    \end{enumerate}
\end{defn}

This definition captures operator-induced exclusions: both goals are individually reachable, but action effects prevent their simultaneous satisfaction. Note that mutex is a \emph{symmetric} relation: if $(g_1, g_2)$ is mutex, then $(g_2, g_1)$ is also mutex.

\begin{defn}[Mutex Goals Count]
    Let $\langle P, O, I, G \rangle$ be a planning problem. The measure $I^{\text{scope}}_{\text{MX}}$ counting goal propositions involved in mutex relationships is defined as:
    \begin{equation*}
        I^{\text{scope}}_{\text{MX}}(\langle P, O, I, G \rangle) = |\{g \in G \mid \exists g' \in G, g \neq g' : \text{mutex}(g, g')\}|
    \end{equation*}
\end{defn}

This measure indicates the \emph{scope} of mutex problems at the goal level: how many goal propositions participate in at least one mutex relationship. The value ranges from $0$ (no mutex) to $|G|$ (all goals involved in mutex).

\begin{defn}[Mutex Pair Count]
    Let $\langle P, O, I, G \rangle$ be a planning problem. The measure $I^{\text{struct}}_{\text{MX}}$ counting unordered mutex goal pairs is defined as:
    \begin{equation*}
        I^{\text{struct}}_{\text{MX}}(\langle P, O, I, G \rangle) = |\{\{g_1, g_2\} \mid g_1, g_2 \in G, g_1 \neq g_2, \text{mutex}(g_1, g_2)\}|
    \end{equation*}
\end{defn}

This measure reveals the \emph{structure} of mutex relationships by counting unordered pairs. Since mutex is symmetric, we count unordered sets $\{g_1, g_2\}$ rather than ordered pairs. The value ranges from $0$ to $\binom{|G|}{2}$. The ratio between $I^{\text{scope}}_{\text{MX}}$ and $I^{\text{struct}}_{\text{MX}}$ reveals conflict patterns: a small ratio suggests isolated mutex pairs, while a large ratio (approaching $\binom{|G|}{2}$) indicates dense mutex cliques. Together, $I^{\text{scope}}_{\text{MX}}$ and $I^{\text{struct}}_{\text{MX}}$ form the Goal Mutex Profile.

\paragraph{Example: Light Switch (Scenario~3).} For the Light Switch scenario defined in Section~\ref{sec:test_scenarios}, forward reachability analysis yields $S_{\text{reach}} = \{\{\textit{off}\}, \{\textit{on}\}\}$. Both goal propositions appear in reachable states, but no state contains both simultaneously. The Goal Mutex Profile yields:
\begin{align*}
    I^{\text{scope}}_{\text{MX}}  & = |\{\textit{on}, \textit{off}\}| = 2     \\
    I^{\text{struct}}_{\text{MX}} & = |\{\{\textit{on}, \textit{off}\}\}| = 1
\end{align*}
The profile detects operator-induced mutual exclusion despite reversibility: actions allow transitions between states, but effects structurally prevent co-occurrence.

\paragraph{Example: Traffic Light (Scenario~4).} For the Traffic Light scenario, reachability analysis yields $S_{\text{reach}} = \{\{\textit{red}\}, \{\textit{green}\}, \{\textit{yellow}\}\}$. All three goal propositions are individually achievable (the system is cyclic), but no state contains more than one color. The Goal Mutex Profile yields:
\begin{align*}
    I^{\text{scope}}_{\text{MX}}  & = |\{\textit{red}, \textit{yellow}, \textit{green}\}| = 3                                                            \\
    I^{\text{struct}}_{\text{MX}} & = |\{\{\textit{red}, \textit{yellow}\}, \{\textit{red}, \textit{green}\}, \{\textit{yellow}, \textit{green}\}\}| = 3
\end{align*}
This demonstrates a complete mutex clique: all goal pairs are mutually exclusive, yielding $I^{\text{struct}}_{\text{MX}} = \binom{3}{2} = 3$. The profile distinguishes this dense conflict structure from isolated mutex pairs.

\subsubsection{Proposal 3: Goal Sequencing Conflict Profile (Formal)}

The Goal Sequencing Conflict Profile quantifies Category~2c inconsistency by analyzing irreversible action-based exclusions where achieving certain goals permanently prevents other goals from being achieved. Unlike mutex relationships (Proposal~2) where goals alternate between states via reversible actions, sequencing conflicts involve one-way transitions creating permanent dead-ends. This profile consists of two complementary measures.

\begin{defn}[Sequencing Conflict]
    \label{defn:sequencing}
    Let $\langle P, O, I, G \rangle$ be a planning problem with reachable states $S_{\text{reach}}$. For goal propositions $g_1, g_2 \in G$ with $g_1 \neq g_2$, a \emph{sequencing conflict} $(g_1, g_2)$ exists if:
    \begin{enumerate}
        \item $\exists s \in S_{\text{reach}} : g_1 \in s$ (goal $g_1$ is achievable)
        \item $\exists s \in S_{\text{reach}} : g_2 \in s$ (goal $g_2$ is achievable)
        \item $\forall s \in S_{\text{reach}}$ with $g_1 \in s$:
              $\nexists s' \in \text{Reach}(s) : g_2 \in s'$ (from every $g_1$-state,
              $g_2$ is unreachable)
    \end{enumerate}
\end{defn}

Informally, $(g_1, g_2)$ is a sequencing conflict when both goals are individually achievable, but achieving $g_1$ creates a permanent dead-end from which $g_2$ can never be satisfied. The requirement that $g_2$ be achievable (condition 2) ensures the conflict is meaningful: it distinguishes ``$g_1$ blocks $g_2$'' from ``$g_2$ was never achievable anyway.''

\textbf{Remark on Irreversibility.} The definition captures \emph{permanent} blocking: once any $g_1$-satisfying state is reached, $g_2$ becomes unreachable. The key semantic property is irreversibility: achieving $g_1$ creates a ``point of no return'' from which $g_2$ can never be satisfied. This contrasts with mutex relationships (Definition~\ref{defn:mutex}), where goals may alternate between states via reversible actions but simply cannot co-occur.

\begin{defn}[Conflicting Goal Propositions Count]
    Let $\langle P, O, I, G \rangle$ be a planning problem. The measure $I^{\text{scope}}_{\text{GS}}$ counting goal propositions involved in sequencing conflicts is defined as:
    \begin{equation*}
        I^{\text{scope}}_{\text{GS}}(\langle P, O, I, G \rangle) = |\{g \in G \mid \exists g' \in G : (g, g') \text{ or } (g', g) \text{ is a sequencing conflict}\}|
    \end{equation*}
\end{defn}

This measure indicates the \emph{scope} of sequencing conflicts: how many goal propositions participate in at least one action-based exclusion. The value ranges from $0$ (no sequencing conflicts) to $|G|$ (all goals involved).

\begin{defn}[Sequencing Conflict Pair Count]
    Let $\langle P, O, I, G \rangle$ be a planning problem. The measure $I^{\text{struct}}_{\text{GS}}$ counting ordered sequencing conflict pairs is defined as:
    \begin{equation*}
        I^{\text{struct}}_{\text{GS}}(\langle P, O, I, G \rangle) = |\{(g_1, g_2) \mid g_1, g_2 \in G, g_1 \neq g_2, (g_1, g_2) \text{ is a sequencing conflict}\}|
    \end{equation*}
\end{defn}

This measure reveals the \emph{structure} of sequencing conflicts by counting ordered pairs. Unlike mutex pairs (which are unordered), sequencing conflicts have direction: $(g_1, g_2)$ means achieving $g_1$ blocks $g_2$, but the reverse may not hold. The value ranges from $0$ to $|G| \times (|G|-1)$. Together, $I^{\text{scope}}_{\text{GS}}$ and $I^{\text{struct}}_{\text{GS}}$ form the Goal Sequencing Conflict Profile.

\paragraph{Example: Trust and Travel (Scenario~5).} For the Trust and Travel scenario defined in Section~\ref{sec:test_scenarios}, forward reachability yields states including $\{\textit{at\_local}, \textit{trust}, \textit{partnership}\}$ and $\{\textit{at\_local}, \textit{opportunity}\}$, but no state contains both goals.

Both goals are individually achievable but never coexist, satisfying the mutex definition. Additionally, $(\textit{opportunity}, \textit{partnership})$ is a sequencing conflict: from any state with \textit{opportunity}, the agent lacks \textit{trust} to achieve \textit{partnership}. However, $(\textit{partnership}, \textit{opportunity})$ is \emph{not} a sequencing conflict: from \textit{partnership}-states, \textit{opportunity} remains reachable via \textit{travel} $\rightarrow$ \textit{exploit} (though this path deletes \textit{partnership}).

The Goal Mutex Profile yields:
\begin{align*}
    I^{\text{scope}}_{\text{MX}}  & = |\{\textit{partnership}, \textit{opportunity}\}| = 2     \\
    I^{\text{struct}}_{\text{MX}} & = |\{\{\textit{partnership}, \textit{opportunity}\}\}| = 1
\end{align*}

The Goal Sequencing Conflict Profile yields:
\begin{align*}
    I^{\text{scope}}_{\text{GS}}  & = |\{\textit{partnership}, \textit{opportunity}\}| = 2   \\
    I^{\text{struct}}_{\text{GS}} & = |\{(\textit{opportunity}, \textit{partnership})\}| = 1
\end{align*}

This demonstrates that irreversible problems exhibit both mutex and sequencing: the goals cannot coexist (mutex), and the sequencing measure additionally identifies the direction of the conflict.

\paragraph{Example: Rival Alliances (Scenario~6).} For the Rival Alliances scenario, forward reachability yields $S_{\text{reach}} = \{\{\textit{neutral\_A}, \textit{neutral\_B}\}, \{\textit{neutral\_A}, \textit{allied\_B}\}, \{\textit{neutral\_B}, \textit{allied\_A}\}\}$. Both goals are individually achievable but never coexist, and allying with either faction alienates the other.

The Goal Mutex Profile yields:
\begin{align*}
    I^{\text{scope}}_{\text{MX}}  & = |\{\textit{allied\_A}, \textit{allied\_B}\}| = 2     \\
    I^{\text{struct}}_{\text{MX}} & = |\{\{\textit{allied\_A}, \textit{allied\_B}\}\}| = 1
\end{align*}

The Goal Sequencing Conflict Profile yields:
\begin{align*}
    I^{\text{scope}}_{\text{GS}}  & = |\{\textit{allied\_A}, \textit{allied\_B}\}| = 2                                             \\
    I^{\text{struct}}_{\text{GS}} & = |\{(\textit{allied\_A}, \textit{allied\_B}), (\textit{allied\_B}, \textit{allied\_A})\}| = 2
\end{align*}

This demonstrates symmetric exclusions where both alliances block each other. Like Trust and Travel, it exhibits both mutex and sequencing, but with $I^{\text{struct}}_{\text{GS}} = 2$ (bidirectional) rather than $1$ (unidirectional). The sequencing structure value thus distinguishes symmetric from asymmetric irreversibility.
